\section{A formalized classical completion input}
The following classical codimension-one case is proved in Lean, without
assuming any Hadamard-completion theorem. It is an input toward the
higher-codimension results. Completion through codimension three is proved
below; the classical codimension-four through seven input remains explicit.
\begin{theorem}[Codimension-one completion]
For every integer $n\ge0$, every real sign matrix $W$ with $n$ rows and
$n+1$ columns satisfying $WW^{\mathsf T}=(n+1)I_n$ extends to a Hadamard
matrix of order $n+1$, retaining every prescribed row.
\end{theorem}
\begin{proof}
Divide each row by $\sqrt{n+1}$ and extend the resulting orthonormal
family to an orthonormal basis of $\mathbb R^{n+1}$. Multiply the basis
rows by $\sqrt{n+1}$ to obtain a square real matrix $H$ with
$HH^{\mathsf T}=(n+1)I$. The square Gram identity also gives
$H^{\mathsf T}H=(n+1)I$. For each column the first $n$ entries are
prescribed signs, so their squares sum to $n$. Its remaining entry
therefore has square one and is a sign. All rows of $H$ are sign rows,
and the required Hadamard identity and row retention hold.
\end{proof}
The exact public declaration is
\texttt{CompletionTheory.codimension\_one\_completion}; its conclusion
uses the standard mathlib \texttt{Matrix.IsHadamard} predicate.

\begin{theorem}[Codimension-two completion]
Every real sign frame of even order with exactly two missing rows extends
to a real Hadamard matrix, retaining all prescribed rows.
\end{theorem}
\begin{proof}
Extend the normalized frame to a real orthonormal basis and rescale, as
above. Let $u_j,v_j$ denote the two missing entries of column $j$.
The column Gram identity gives $u_j^2+v_j^2=2$. Fix a column $j_0$,
put $a=u_{j_0}$, $b=v_{j_0}$, and rotate the missing rows by
$\frac12\left(\begin{smallmatrix}a+b&b-a\\a-b&a+b\end{smallmatrix}\right)$.
This is orthogonal and sends their $j_0$ column to $(1,1)^T$.
Every residual correlation with $j_0$ is an even integer, since the
order and number of prescribed signs are even. Consequently each
new missing column satisfies $x^2+y^2=2$ and $x+y\in2\mathbb Z$.
The latter sum is $-2,0$, or $2$ by Cauchy--Schwarz, and each case
forces $x,y\in\{1,-1\}$. Thus the rotated missing rows are sign rows.
\end{proof}
The declaration is \texttt{CompletionTheory.codimension\_two\_completion}.

\begin{theorem}[A switched codimension-three criterion]
Let $P$ be a real partial Hadamard matrix with exactly three missing
rows. If every off-diagonal column inner product of $P$ belongs to
$\{-3,1\}$, then $P$ extends to a real Hadamard matrix retaining all rows.
\end{theorem}
\begin{proof}
Extend to a scaled real orthonormal basis. The three residual coordinates
$r_j$ of each column have squared norm three and pairwise inner product
three or minus one. Inner product three implies equality of the residual
vectors. If $t$ distinct residual vectors occur, their sum has squared
norm $3t-t(t-1)=t(4-t)$, so $t\le4$. Assign the distinct classes to distinct
members of $(1,1,1),(1,-1,-1),(-1,1,-1),(-1,-1,1)$. These sign vectors have
exactly the same Gram matrix as the residual columns. The resulting
three sign rows, stacked with $P$, have column Gram equal to the order
times identity. The square matrix therefore also has the required row
Gram and is Hadamard.
\end{proof}
This is a sufficient criterion with an explicit correlation hypothesis,
not the generic codimension-three completion theorem. Its declaration in
the \texttt{CompletionTheory} namespace is
\texttt{tetrahedral\_completion\_of\_card}; the separate
\texttt{tetrahedral\_sign\_factor} proves the tetrahedral factorization.

\begin{theorem}[Classical completion through codimension three]
Every partial real Hadamard matrix of positive order divisible by four
with at most three missing rows extends to a Hadamard matrix retaining
all its prescribed rows.
\end{theorem}
\begin{proof}
The cases zero, one and two were handled above. For three missing rows,
the number $r$ of prescribed rows is congruent to one modulo four.
Multiply each column by the product of its prescribed entries. Since
$r$ is odd, every switched column has product one. The entrywise
product of any two switched columns also has product one, hence has an
even number of negative signs. Its sum is consequently congruent to
$r$, and thus to one, modulo four. The real residual columns have
squared norm three, so Cauchy--Schwarz bounds every off-diagonal
prescribed column correlation between minus three and three. Only
minus three and one have the required congruence. Apply the preceding
tetrahedral criterion and undo the column switches. All prescribed rows
are retained. Every step, including the parity argument, is formalized.
\end{proof}
The arbitrary-index theorem
\texttt{CompletionTheory.codimension\_le\_three\_completion}
returns the standard mathlib Hadamard predicate and an injection identifying
every original row. This proves only the stated classical small-defect
case; the codimension-four through seven inputs remain unformalized.
